\scanpage{112}
\mspage{99}
% archive sheet number: 111

b) Soit $I_{0}$ la catégorie discrète ayant $=$
$=$ ensemble d'objets que $I$, et
\[
 \unsure{v|u}_{I}^{\deleted{\unsure{*|I}}}\ \text{ou }v:I_{0}\mathrel{\unsure{-|\longrightarrow}}I\tag*{(6.1.4)}
\]
le foncteur canonique. \unsure{Distinguons|Désignons} par
\deleted{\[
 C(\unsure{I|\,})\tag*{(6.1.5)}
\]}
\[
 C(I)=(TF(I))^{*},\qquad TC(I)=C(I)\cap W(I)=F(I)^{*}\tag*{(6.1.5)}
\]
les ensembles de cofibr. resp. de cofibr.
triviales dans $M(I)$, et \unsure{utilisons} les
$=$s notations $C(I_{0}),TC(I_{0})$ pour \unsure{Do.|Do}. On
a alors
\[
\begin{aligned}
 &\deleted{\unsure{C(I_{0})}}\ u:X_{0}\mathrel{-}Y_{0}\ \text{dans }M(I_{0})\\
 &\text{est dans } C_{0}\ \text{(resp. }TC_{0})\ \text{ssi pour}\\
 &\text{tt }i\ \text{dans }\unsure{i|I},\quad u(i):X_{0}(i)\mathrel{-}Y_{0}(i)\\
 &\text{est dans }C\ \text{(resp. }TC).
\end{aligned}\tag*{(6.1.6)}
\]
Enfin, on a
\[
\left\{\begin{aligned}
 C(I)&=(v_{!}(C(I_{0})))^{\sim},\\
 TC(I)&=(v_{!}(TC(I_{0})))^{\sim}.
\end{aligned}\right.\tag*{(6.1.7.)}
\]

c) Soit \deleted{\unsure{un}}
\[
 u:I\longrightarrow I'
\]
une flèche \deleted{\unsure{quel}} dans Cat, considérons
le couple de foncteurs adjoints
\diagram{\begin{tikzcd}[ampersand replacement=\&,column sep=large]
{(6.1.8)} \& M(I) \arrow[r,shift left=1.2ex,"{u_{!}}"] \& M(I') \arrow[l,shift left=1.2ex,"{u^{*}}"]
\end{tikzcd}}.

